\chapter{Conception de systèmes}\label{ch:iv3:system-design}

\begin{ivtrains}
Concevoir des systèmes de trading à voix haute~: moteurs d'appariement (matching engine), feed handlers, contrôles de risque, bases de ticks (tick stores), routeurs et plateformes de recherche. Partir d'un chiffre, nommer la propriété qui ne doit jamais céder, et dire ce qu'on accepterait de sacrifier. Le code présenté ici est constitué de petits modèles Python de ces propriétés, et chaque énoncé précise ses hypothèses. Les chiffres donnés dans un énoncé sont des hypothèses de l'exercice, et non des mesures d'un système quelconque.
\end{ivtrains}

\begin{question}[$\star$ \iqroles{QD, Dev, LLE} \iqfirm{toutes} \iqlevel{junior}]\label{iq:iv3:system-design:1}
Un feed handler reçoit 50\,000 messages par seconde, et un message passe en moyenne 40 microsecondes à l'intérieur, attente comprise. Combien de messages se trouvent-ils à l'intérieur à un instant donné, en moyenne~? Plus tard, vous observez une profondeur stable de 200 au même débit. Quel est alors le temps moyen passé à l'intérieur~? Quelle est la loi que vous avez utilisée, et que suppose-t-elle~?
\end{question}

\begin{question}[$\star$ \iqroles{QD, Dev, LLE} \iqfirm{toutes} \iqlevel{junior}]\label{iq:iv3:system-design:2}
Une requête appelle 10 backends en parallèle et attend la réponse de chacun d'eux. Chaque backend, indépendamment, dépasse 5 ms sur 1~\% des appels. Quelle fraction des requêtes dépasse 5 ms~? Et avec 100 backends~? Pourquoi cela fait-il du 99\up{e}~centile le chiffre à suivre, et non la moyenne~?
\end{question}

\begin{question}[$\star$ \iqroles{QD, Dev, Analyste} \iqfirm{toutes} \iqlevel{junior}]\label{iq:iv3:system-design:3}
Un feed atteint un pic de 1 million de messages par seconde et une moyenne de 200\,000 sur une session de 6.5~heures, avec 40 octets par message. En négligeant les en-têtes, quel débit de liaison le pic exige-t-il~? Quel volume de données brutes représente une journée, et quel volume une année de 252 jours~? Pour lequel de ces deux débits dimensionnez-vous chaque partie du système~?
\end{question}

\begin{question}[$\star$ \iqroles{QD, Dev} \iqfirm{toutes} \iqlevel{junior}]\label{iq:iv3:system-design:4}
Un service est disponible 99.9~\% du temps. Combien d'heures d'indisponibilité cela représente-t-il sur une année~? Quelle est la disponibilité de deux services de ce type lorsqu'une requête a besoin des deux, puis lorsqu'un seul des deux suffit, en supposant qu'ils tombent en panne de façon indépendante~? Pourquoi cette dernière hypothèse est-elle le point faible~?
\end{question}

\begin{question}[$\star$ \iqroles{QD, Dev} \iqfirm{toutes} \iqlevel{junior}]\label{iq:iv3:system-design:5}
Un client envoie un ordre, dépasse le délai d'attente, puis le renvoie. En réalité, le premier est bien arrivé, et seule la réponse s'est perdue. Quelles sont les deux lignes imprimées par le \cref{lst:iv3:sd:retry}~? À quoi sert l'identifiant d'ordre du client, et pendant combien de temps la passerelle doit-elle en garder la trace~? Que doit-il être vrai des deux écritures de la seconde branche si la passerelle peut tomber en panne entre les deux~?
\omcode{code/interviews-3/06-system-design/python/retry.py}{1}{22}{Une nouvelle tentative après une réponse perdue.}\label{lst:iv3:sd:retry}
\end{question}

\begin{question}[$\star$ \iqroles{QD, Dev, LLE} \iqfirm{toutes} \iqlevel{junior}]\label{iq:iv3:system-design:6}
Un processus de trading effectue ces opérations pour chaque mise à jour des données de marché~: décider s'il faut envoyer un ordre, vérifier les limites de risque, envoyer l'ordre, mettre à jour la position, écrire un enregistrement d'audit dans une base de données, formater une ligne de journal et rafraîchir un tableau de bord. Lesquelles relèvent du chemin critique (hot path), et lesquelles non~? Comment les autres obtiennent-elles leurs données~?
\end{question}

\begin{question}[$\star$ \iqroles{QD, Dev} \iqfirm{toutes} \iqlevel{junior}]\label{iq:iv3:system-design:7}
10\,000 clients veulent chacun un seul cours. Comparez le polling, une fois par seconde, à un serveur qui pousse chaque changement. Combien de messages par seconde le serveur envoie-t-il dans chaque conception, si le cours change deux fois par seconde, et s'il change une fois par minute~? Quelle est l'ancienneté moyenne de la donnée avec le polling~? Dans quel cas le polling est-il le meilleur choix~?
\end{question}

\begin{question}[$\star$ \iqroles{QD, Dev, LLE} \iqfirm{toutes} \iqlevel{junior}]\label{iq:iv3:system-design:8}
Un seul thread traite un ordre en 2 microsecondes. Quel est son débit maximal~? Si les ordres arrivent au hasard à 400\,000 par seconde et que le temps de service suit une loi exponentielle, quel est le temps moyen passé dans le système~? Pourquoi ne peut-on pas prévoir de le faire tourner à 100~\%~?
\end{question}

\begin{question}[$\star\star$ \iqroles{QD, Dev} \iqfirm{toutes} \iqlevel{junior, confirmé}]\label{iq:iv3:system-design:9}
Un cache de symboles est réparti sur 10 serveurs avec \texttt{hash(key) \% n}. Vous en ajoutez un onzième. Quelle fraction des clés change-t-elle de serveur~? En Python, écrivez un anneau de hachage avec des nœuds virtuels, où une clé va au point suivant dans le sens horaire. Quelle fraction des clés se déplace alors, et vers où~? À quoi servent les nœuds virtuels~? Quand le sharding par hachage du symbole est-il un mauvais choix en trading~?
\end{question}

\begin{question}[$\star\star$ \iqroles{QD, Dev} \iqfirm{toutes} \iqlevel{junior, confirmé}]\label{iq:iv3:system-design:10}
En Python, écrivez le cycle de vie d'un ordre tel que l'émetteur le voit~: new, open, pending cancel, et les trois états finaux filled, cancelled et rejected. Une annulation peut croiser la dernière exécution sur le fil, et une plateforme peut annuler un ordre sans que personne le lui ait demandé. Que doit faire le système dans chacun de ces cas~? Quels événements doivent déclencher une alerte~? Pourquoi ne pas traiter tout événement inattendu comme anodin~?
\end{question}

\begin{question}[$\star\star$ \iqroles{QD, LLE, Risque} \iqfirm{toutes} \iqlevel{junior, confirmé}]\label{iq:iv3:system-design:11}
Écrivez, en Python, un contrôle pré-trade avec une limite de taille, un tunnel de prix (price collar) autour d'un prix de référence, et une limite de position. Les ordres ouverts comptent~: la limite de position doit tenir même si tous les ordres ouverts sont ensuite exécutés. Pourquoi un contrôle qui ne porte que sur la position actuelle est-il erroné~? Où doivent se situer ces contrôles, et quel coût doivent-ils avoir sur le chemin critique~?
\end{question}

\begin{question}[$\star\star$ \iqroles{QD, Dev, Risque} \iqfirm{toutes} \iqlevel{confirmé}]\label{iq:iv3:system-design:12}
Concevez un kill switch pour un système de trading. Que doit-il faire, dans quel ordre, et qui ou quoi peut le déclencher~? Pourquoi ne doit-il pas dépendre du processus qu'il arrête~? De quoi ne vous protégera-t-il pas~? Quels déclencheurs automatiques mettriez-vous en place~?
\end{question}

\begin{question}[$\star\star$ \iqroles{QD, LLE} \iqfirm{crypto, teneur de marché, hedge fund} \iqlevel{junior, confirmé}]\label{iq:iv3:system-design:13}
Une plateforme publie des mises à jour numérotées du carnet, chacune donnant un niveau de prix et sa nouvelle taille, et fournit séparément un instantané étiqueté par le numéro de séquence qu'il reflète. En Python, construisez le carnet à partir de ces données. Vous rejoignez le flux en cours, si bien qu'un instantané peut arriver avant ou après les mises à jour qui l'entourent, et que des mises à jour peuvent être perdues. Que faites-vous des mises à jour reçues avant l'arrivée d'un instantané~? Que se passe-t-il si l'instantané est plus ancien que la plus ancienne de vos mises à jour tamponnées~? Que garantit la conception~?
\end{question}

\begin{question}[$\star\star$ \iqroles{QD, Dev, LLE} \iqfirm{toutes} \iqlevel{junior, confirmé}]\label{iq:iv3:system-design:14}
Un consommateur ne parvient pas à suivre les mises à jour de prix de 10 symboles. En Python, écrivez un tampon de conflation~: \texttt{put(symbol, value)} et \texttt{take(n)}, qui ne garde que la valeur la plus récente de chaque symbole, de façon qu'un symbole très actif ne puisse pas priver un symbole calme de ses mises à jour. Quelle quantité de mémoire peut-il utiliser~? Quel type de données ne doit jamais faire l'objet d'une conflation~?
\end{question}

\begin{question}[$\star\star$ \iqroles{QD, Dev} \iqfirm{toutes} \iqlevel{junior, confirmé}]\label{iq:iv3:system-design:15}
Un publisher envoie 10\,000 messages par seconde de 100 octets chacun, sur TCP, avec un tampon de 16 MiB pour chaque client. Un client lit à la moitié de ce débit. Au bout de combien de temps son tampon est-il plein~? Quelles sont alors les options~? Lesquelles retenez-vous pour un tableau de bord, pour un système de risque et pour un moteur de trading qui s'abonnent tous~?
\end{question}

\begin{question}[$\star\star$ \iqroles{QD, Dev} \iqfirm{crypto} \iqlevel{junior, confirmé}]\label{iq:iv3:system-design:16}
Concevez une page web de cours en temps réel pour 100 instruments, avec 1 million de clients connectés. Chaque cours change une fois par seconde, et une mise à jour pèse 50 octets. Quel débit sortant cela exige-t-il~? Et si chaque client ne suit que 10 cours~? Si un serveur tient 20\,000 connexions, combien de serveurs faut-il~? Comment les mises à jour leur parviennent-elles, et que se passe-t-il quand un serveur tombe en panne~?
\end{question}

\begin{question}[$\star\star$ \iqroles{QD, LLE, QT} \iqfirm{crypto, teneur de marché, hedge fund} \iqlevel{confirmé}]\label{iq:iv3:system-design:17}
Deux clients colocalisés sur une même bourse comparent leurs journaux, et ils affichent un meilleur cours acheteur (best bid) différent au même instant d'horloge murale. Est-ce un bogue~? Dressez la liste de ce qui pourrait l'expliquer. Comment établiriez-vous laquelle de ces causes est en jeu, et qu'auriez-vous dû enregistrer~?
\end{question}

\begin{question}[$\star\star$ \iqroles{QD, DS, QR} \iqfirm{toutes} \iqlevel{junior, confirmé}]\label{iq:iv3:system-design:18}
Le \Cref{lst:iv3:sd:columns} écrit quatre jours de ticks au format Parquet, dans un répertoire par date, et interroge une colonne d'un jour avec DuckDB 1.5. Qu'affiche-t-il, à peu près, et que signifient ses contrôles~? Pourquoi un format en colonnes convient-il à cette requête, et quand ne convient-il pas~? Comment organiseriez-vous une base de ticks pour une année~? Pourquoi chaque fichier ne correspond-il pas à un seul symbole~?
\omcode{code/interviews-3/06-system-design/python/iv_columns.py}{2}{31}{Une base de ticks en miniature.}\label{lst:iv3:sd:columns}
\end{question}

\begin{question}[$\star\star$ \iqroles{QD, DS, LLE} \iqfirm{toutes} \iqlevel{junior}]\label{iq:iv3:system-design:19}
Les horodatages sont exprimés en nanosecondes sur 64 bits, avec un tick par microseconde. En Python, écrivez un codec~: différences entre voisins, codage zigzag pour les valeurs négatives, puis entiers à longueur variable. Combien d'octets occupent 10\,000 horodatages de ce type~? Pourquoi les deltas sont-ils si petits~? Qu'advient-il d'une suite de valeurs aléatoires non triées~? Pourquoi une base de ticks commence-t-elle par cette étape avant un compresseur général~?
\end{question}

\begin{question}[$\star\star$ \iqroles{QD, QR, DS} \iqfirm{toutes} \iqlevel{junior, confirmé}]\label{iq:iv3:system-design:20}
Un fournisseur corrige un prix un jour après l'avoir publié. En Python, écrivez une table avec un instant d'événement et un instant de connaissance, ainsi que \texttt{as\_of(known\_at)}, qui renvoie la table telle qu'elle était alors. Pourquoi ne pas simplement écraser la ligne~? Que change l'écrasement pour un backtest, et pour un rapport que vous avez envoyé le mois dernier~?
\end{question}

\begin{question}[$\star\star$ \iqroles{QD, QT, LLE} \iqfirm{toutes} \iqlevel{junior, confirmé}]\label{iq:iv3:system-design:21}
Vous devez acheter 300 actions. La plateforme A facture des frais de 0.003 par action et affiche 100 à 10.00 et 200 à 10.01. La plateforme B facture des frais de 0.001, avec 150 à 10.00 et 300 à 10.02. La plateforme C verse un rebate de 0.002, avec 100 à 10.01 et 500 à 10.03. En Python, répartissez l'ordre selon le prix net. Quel est le plan, et quel en est le coût~? Que diriez-vous de cette réponse, sachant que ces ordres mettent du temps à arriver~?
\end{question}

\begin{question}[$\star\star$ \iqroles{QD, Dev, Risque} \iqfirm{toutes} \iqlevel{junior, confirmé}]\label{iq:iv3:system-design:22}
Chaque jour, vous comparez votre propre registre des exécutions avec la copie distincte de la plateforme. En Python, écrivez la comparaison, indexée par identifiant d'exécution. Que signifie chacun des trois types d'écarts~? Lequel est le plus urgent, et pourquoi~? Pourquoi lancer ce contrôle en cours de journée, en plus de la fin de journée~?
\end{question}

\begin{question}[$\star\star$ \iqroles{QD, Dev} \iqfirm{crypto, banque} \iqlevel{junior, confirmé}]\label{iq:iv3:system-design:23}
Concevez un grand livre des comptes, et implémentez un virement en Python sur SQLite. Il ne doit jamais perdre ni créer d'argent, ne jamais mettre un compte à découvert, survivre à un arrêt brutal en cours de traitement, et traiter une requête relancée comme un seul virement. Pourquoi stocker des écritures plutôt qu'un solde~? À quoi servirait une colonne de solde~? Où doivent se trouver la vérification et les écritures, pour la concurrence~?
\end{question}

\begin{question}[$\star\star$ \iqroles{QD, LLE} \iqfirm{crypto, teneur de marché, hedge fund} \iqlevel{junior, confirmé}]\label{iq:iv3:system-design:24}
Un message traverse trois étapes successives, et chacune dure 10 microsecondes, sauf que, avec une probabilité de 2~\%, indépendamment, elle dure 100 microsecondes. Quel est le 99\up{e}~centile d'une étape~? Et celui du total~? Peut-on alors additionner les centiles dans un budget de latence~? En Python, calculez exactement la distribution du total. Quand la somme des centiles est-elle finalement juste~?
\end{question}

\begin{question}[$\star\star$ \iqroles{QD, LLE} \iqfirm{toutes} \iqlevel{junior, confirmé}]\label{iq:iv3:system-design:25}
Un serveur unique a un temps de service moyen de 1 microseconde, avec des arrivées aléatoires et un service exponentiel. Quel est le temps moyen passé dans le système à 50~\%, 90~\% et 99~\% de charge~? Vérifiez ces valeurs par simulation en Python, avec une ligne par client. Que change un temps de service fixe~? Qu'est-ce que cela implique pour la façon dont vous faites tourner un thread critique~? Que masque la moyenne~?
\end{question}

\begin{question}[$\star\star$ \iqroles{QD, LLE} \iqfirm{toutes} \iqlevel{confirmé}]\label{iq:iv3:system-design:26}
Un journal doit être écrit sur disque avant qu'un ordre ne soit acquitté, et un flush prend 2 ms sur ce disque, ce qui est une hypothèse de l'exercice. Quel est le débit maximal avec un flush par ordre~? Qu'est-ce que le group commit~? Simulez-le en Python à 10\,000 ordres par seconde~: quelle est la taille des lots, et quelle est la latence moyenne~? Que se passe-t-il à 300 ordres par seconde~? Que changerait-il si l'on répliquait au lieu de faire un flush~?
\end{question}

\begin{question}[$\star\star$ \iqroles{QD, QR, DS} \iqfirm{toutes} \iqlevel{junior}]\label{iq:iv3:system-design:27}
Un centre de recherche met en cache les résultats de backtests. En Python, écrivez la clé de cache. Qu'est-ce qui doit y figurer pour que la même clé désigne toujours le même résultat~? Quelles sont les erreurs courantes qui conduisent à se tromper~? Que ne peut une clé capturer, quelle que soit sa construction~?
\end{question}

\begin{question}[$\star\star$ \iqroles{QD, QR, DS} \iqfirm{hedge fund} \iqlevel{junior, confirmé}]\label{iq:iv3:system-design:28}
Un chercheur relance un backtest vieux de six mois avec le même code et obtient un ratio de Sharpe différent. Listez les causes, de la plus probable à la moins probable. Comment concevriez-vous une plateforme de recherche permettant de reproduire un résultat, ou d'expliquer l'écart~?
\end{question}

\begin{question}[$\star\star\star$ \iqroles{QD, LLE, Dev} \iqfirm{crypto, teneur de marché, hedge fund} \iqlevel{confirmé}]\label{iq:iv3:system-design:29}
Un routeur envoie des ordres vers des dizaines de plateformes, et l'une d'elles se bloque. En Python, simulez deux conceptions~: un seul worker pour tout, et un worker par plateforme. Avec 20 ordres, un par seconde, en alternant entre une plateforme saine à 10 ms et une plateforme bloquée à 50 s, en commençant par la saine, quelle est la pire latence de la plateforme saine dans chacun des deux cas~? Au-delà de la file, que doit faire le routeur des ordres en vol vers la plateforme bloquée~? Dans quel état se trouve un tel ordre~?
\end{question}

\begin{question}[$\star\star\star$ \iqroles{QD, LLE, Dev} \iqfirm{toutes} \iqlevel{confirmé}]\label{iq:iv3:system-design:30}
Écrivez en Python un journal en ajout seul~: \texttt{append(f, payload)}, et \texttt{recover(data)}, qui renvoie le plus long préfixe correct après une panne, quelle que soit la façon dont le fichier a été tronqué. Quel est le rôle de chaque champ de l'enregistrement~? Comment le testez-vous~? Si une journée compte 5 milliards d'événements et qu'un rejeu tourne à un million par seconde, combien de temps dure la reprise, et que faites-vous pour la réduire~?
\end{question}

\begin{question}[$\star\star\star$ \iqroles{QD, LLE} \iqfirm{crypto, teneur de marché} \iqlevel{confirmé}]\label{iq:iv3:system-design:31}
Beaucoup de moteurs d'appariement font tourner le carnet sur un seul thread. Pourquoi~? À 2 microsecondes par ordre, qu'est-ce qui limite cette conception, et comment dépasser cette limite~? Quel est le problème d'un ordre qui touche deux partitions~? Pourquoi cette conception est-elle aussi celle qui rend la réplication simple~?
\end{question}

\begin{question}[$\star\star\star$ \iqroles{QD, Dev} \iqfirm{toutes} \iqlevel{confirmé}]\label{iq:iv3:system-design:32}
Un consommateur lit un journal durable, met à jour un état et enregistre jusqu'où il a lu, et il peut tomber en panne à chaque étape. En Python, modélisez trois conceptions~: appliquer puis enregistrer, enregistrer puis appliquer, et les deux dans une seule étape atomique. Laquelle surcompte, laquelle sous-compte, et laquelle est correcte, avec un test qui permette de le montrer~? Que signifie réellement l'exactly-once ici~? Et si l'état se trouve sur un autre système qui n'a pas de transactions avec le décalage~?
\end{question}

\begin{question}[$\star\star\star$ \iqroles{QD, Dev, LLE} \iqfirm{toutes} \iqlevel{confirmé}]\label{iq:iv3:system-design:33}
Un primaire se met en pause pendant longtemps, il est déclaré mort, et un standby prend la relève. Puis l'ancien se réveille, persuadé qu'il est toujours le primaire. Quelles sont les deux lignes qu'imprime le \cref{lst:iv3:sd:fencing}~? Que fait l'époque, et qui doit la vérifier~? Pourquoi un délai d'expiration du heartbeat ne suffit-il pas à lui seul~? Que se passe-t-il si l'effet de l'ancien primaire est un ordre envoyé à une plateforme, et non une écriture dans ce magasin~?
\omcode{code/interviews-3/06-system-design/python/iv_fencing.py}{1}{23}{Protection d'un primaire revenu à la vie.}\label{lst:iv3:sd:fencing}
\end{question}

\begin{question}[$\star\star\star$ \iqroles{QD, Dev} \iqfirm{toutes} \iqlevel{confirmé}]\label{iq:iv3:system-design:34}
L'état est répliqué de façon synchrone sur trois sites, avec des allers-retours de 0 ms (l'écrivain lui-même), 1 ms et 40 ms, hypothèses de l'exercice. Quelle est la latence de validation si l'on attend tous les sites, et si l'on attend une majorité~? Combien de pannes chaque conception supporte-t-elle~? Que devient la latence lorsque le site à 1 ms tombe~? Pourquoi les deux majorités doivent-elles se recouvrir, et que casse-t-on si l'on se contente de moins~? Montrez-le en Python.
\end{question}

\begin{question}[$\star\star\star$ \iqroles{QD, QR, DS} \iqfirm{toutes} \iqlevel{confirmé}]\label{iq:iv3:system-design:35}
Un job de recherche tourne pendant des jours sur une machine qui tombe en panne en moyenne une fois toutes les 24 heures. Un point de contrôle coûte 60 s. En notant $C$ le coût d'un point de contrôle, $M$ le temps moyen entre pannes et $T$ l'intervalle, donnez une formule au premier ordre pour la fraction de temps perdue, et trouvez le meilleur $T$. Vérifiez-le par simulation en Python, à l'intervalle optimal et de part et d'autre de celui-ci. Quand vaut-il mieux ne pas s'en soucier~?
\end{question}

\begin{question}[$\star\star\star$ \iqroles{QD, LLE, Dev} \iqfirm{toutes} \iqlevel{confirmé}]\label{iq:iv3:system-design:36}
Une session d'ordres numérote ses messages, et à la reconnexion, l'autre côté demande un renvoi à partir du premier numéro manqué. En Python, écrivez un émetteur et un récepteur. Pourquoi un ordre vieux de plusieurs minutes ne doit-il pas être renvoyé~? Qu'est-ce qui le remplace, et pourquoi ne peut-on pas simplement sauter le numéro~? Après cela, comment savoir ce qu'il est advenu des ordres que vous n'avez pas renvoyés~?
\end{question}

\begin{question}[$\star\star\star$ \iqroles{QD, Dev, Risque} \iqfirm{toutes} \iqlevel{confirmé}]\label{iq:iv3:system-design:37}
Vous devez déployer une nouvelle version d'un système de routage d'ordres réparti sur plusieurs serveurs. Décrivez le déploiement, en précisant ce que vous vérifiez avant l'ouverture, comment on active une nouvelle fonctionnalité, et comment on revient en arrière. Quel est le risque de réutiliser un ancien indicateur de configuration~? Que doit-on garantir pour une version qui tourne à côté de l'ancienne~?
\end{question}

\begin{question}[$\star\star\star\star$ \iqroles{QD, LLE, Dev} \iqfirm{crypto, teneur de marché} \iqlevel{confirmé}]\label{iq:iv3:system-design:38}
Concevez le chemin des ordres d'une bourse pour 100\,000 ordres par seconde sur un groupe d'instruments, depuis la passerelle jusqu'aux données de marché. Nommez chaque composant, ce qu'il stocke et l'ordre des événements. Où vivent les contrôles de risque~? Comment définit-on l'équité~? Comment un standby reste-t-il exact, et qu'a besoin un rejeu pour donner le même résultat~? Votre moteur tombe en panne à un instant connu~: comment reprendre, et comment savoir que rien n'a été perdu~? Écrivez le cœur déterministe en Python.
\end{question}

\begin{question}[$\star\star\star\star$ \iqroles{QD, LLE, Dev} \iqfirm{crypto, teneur de marché, hedge fund} \iqlevel{confirmé}]\label{iq:iv3:system-design:39}
Concevez le service interne de données de marché d'une société de trading. Il reçoit les flux de nombreuses plateformes, les normalise et dessert 300 consommateurs~: moteurs de trading, risque, capture pour la recherche et tableaux de bord. Le pic est de 2 millions de messages par seconde à 64 octets, et la moyenne sur une séance de 6.5 heures est de 500\,000. Dimensionnez-le. Comment les données parviennent-elles à chaque type de consommateur~? Comment détecte-t-on une perte et comment la répare-t-on~? Comment la bascule se fait-elle, et que reste-t-il en service lorsqu'une ligne de plateforme tombe~? Que mesureriez-vous~?
\end{question}

\begin{question}[$\star\star\star\star$ \iqroles{QD, QR, DS} \iqfirm{hedge fund} \iqlevel{confirmé}]\label{iq:iv3:system-design:40}
Concevez une plateforme de recherche pour 50 chercheurs. Les entrées sont 10 ans de ticks et de données de référence, soit environ 5 Gio par jour et par région après compression, pour 25 régions. La simulation d'une journée de stratégie occupe un cœur pendant environ 8 minutes, et une étude typique représente 500 journées de stratégie fois 3 variantes, sur un pool de 64 cœurs. Dimensionnez le stockage et le temps de calcul. Quelles sont les couches~? Comment gère-t-on les versions des données, l'exactitude point-in-time, les coûts et les résultats~? Qu'est-ce qui empêche qu'un bon backtest soit le fruit du hasard~?
\end{question}
